% Source: https://artofproblemsolving.com/community/c3h3047185_mathcounts_state_27
% Stable identity: aops-topic-3047185-q01
The angles at $ P $ and $ R $ subtend the diameter and are right angles. Thus $ QS^2=2^2+6^2=40 $. Triangle $ QRS $ is right isosceles, so its two legs have squared length 20 and area 10. Triangle $ PQS $ has area $\dfrac{1}{2}\cdot2\cdot6=6 $. Adding gives $\boxed{16}$.
\begin{asy}
usepackage("amsmath");
size(7cm);
pair Q=(0,0),PtS=(sqrt(40),0),P=(4/sqrt(40),sqrt(3.6)),R=(sqrt(10),-sqrt(10));draw(circle((Q+PtS)/2,sqrt(10)));draw(P--Q--R--PtS--cycle^^Q--PtS);label("$P$",P,N);label("$Q$",Q,W);label("$R$",R,S);label("$S$",PtS,E);label("$2$",P/2,NW);label("$6$",(P+PtS)/2,NE);
\end{asy}
